\section{附注}\label{sec:notes-on-curvature}

\subsection{第 \ref{sec:manifolds-of-constant-curvature} 节}\label{subsec:notes-manifolds-of-constant-curvature}

双曲平面的等距映射是\textbf{线性分式变换} (或 \textbf{Möbius 变换}) 的例子, 即由
\[
  f(z)=\frac{az+b}{cz+d}
\]
给出的映射 $f:\mathbb{C}\to\mathbb{C}$, 其中 $a,b,c,d\in\mathbb{C}$ 满足 $ad-bc\ne0$. 容易看出, 这些变换中的每一个都是下列类型的变换的复合:
\begin{enumerate}[label=(\arabic*),leftmargin=*]
\item 平移: $z\mapsto z+b$;
\item 旋转: $z\mapsto az$, $|a|=1$;
\item 位似: $z\mapsto rz$, $r>0$;
\item 反演: $z\mapsto1/z$,
\end{enumerate}
从而显然它们把直线和圆映成直线或圆.

可以分别用 $z\to\infty$ 和 $z\to-d/c$ 的极限引入特殊值 $f(\infty)=\frac ac$ 和 $f(-\frac dc)=\infty$. 因此利用球极投影, 可以将 $f$ 看成球面到自身的映射. 注意, 平面上的直线和圆都对应于球面上的圆, 因此可以说, 视为球面上一个映射的 Möbius 变换将圆映成圆.

\subsection{文献附注}\label{subsec:curvature-literature}

本章材料见于大多数 Riemann 几何著作 (例如 [\RGref{bib:ch04-Boo03}, \RGref{bib:ch04-dC93}, \RGref{bib:ch04-GHL04}]). Gauss–Bonnet 定理的证明 (出自 S. Chern) 紧随 [\RGref{bib:ch04-dC93}, \RGref{bib:CCL00}] 的做法. 此定理如何纳入纤维丛示性类的一般理论, 参见 [\RGref{bib:ch04-KN96}, \RGref{bib:Jos02}]. 关于 $\mathbb{R}^3$ 中曲面等距浸入的更初等讨论 (包括 Gauss–Bonnet 定理的证明) 可见 [\RGref{bib:dC76}, \RGref{bib:Mor98}].
